\section{Tensors, Flat Vectors, And Inner Products}

A tensor is a multi-dimensional array.  In this paper the most common tensor is
an order-$\ell$ tensor of shape $2\times\cdots\times2$, meaning it has one field
value for every Boolean address $b=(b_0,\ldots,b_{\ell-1})\in\B^\ell$.  We write
that entry as $x[b]$.

The tensor-first convention is this paper's notation, but the multilinear
extension object is the same one used by HyperPlonk and Blaze: a Boolean table
has a unique multilinear extension to field points~[CBBZ23,Blaze24].

Formally, such a tensor is a function
\[
  x:\B^\ell\to\F.
\]
We write the vector space of all such tensors as
\[
  \F^{\B^\ell}.
\]
Thus saying $x\in\F^{\B^\ell}$ means exactly that $x$ assigns one field element
to each Boolean address in $\B^\ell$.

This tensor notation is the mathematical object.  Implementations eventually
serialize the same entries into a flat vector for memory, Merkle leaves, and
codeword coordinates:
\[
  x=(x[0],x[1],\ldots,x[N-1]).
\]
Here $N=2^\ell$.  The serialization is a deterministic address order for the
Boolean entries.  In this paper it is always the natural binary order.  For
\(0\le i<2^\ell\), write
\[
  \bits_\ell(i)\in\B^\ell
\]
for the length-\(\ell\) binary expansion of \(i\), with leading zeros included.
Then the \(i\)-th serialized entry is the tensor entry
\[
  x[\bits_\ell(i)].
\]
For an integer interval \([s,t]\subseteq[0,2^\ell-1]\), write
\[
  \bits_\ell([s,t])
  :=
  \{\bits_\ell(i):s\le i\le t\}\subseteq\B^\ell.
\]
This notation is only an implementation shorthand for a set of Boolean
addresses.  Protocol statements still use the binary address set.  For product
domains such as \(\B^a\times\B^b\), the natural layout is the lexicographic
binary layout induced by concatenation \((u,v)\mapsto u\Vert v\in\B^{a+b}\).
That order matters for implementors, because it tells them which memory
locations are paired in a fold or adjacent in an accumulator, but it is not the
algebraic object being folded.

The inner product of two tensors of the same shape is the sum of entrywise
products:
\[
  \ip{x}{y}=\sum_{b\in\B^\ell} x[b]y[b].
\]
After serializing both tensors with the same address order, this is implemented
as the usual inner product of two length-$N$ flat vectors:
\[
  \ip{x}{y}=\sum_{i=0}^{N-1} x[i]y[i].
\]
Many claims in this protocol reduce to inner products.  In particular, the
evaluation of a multilinear extension can be written as an inner product between
the original tensor and a tensor of public evaluation weights.  The protocol
uses both views constantly: the tensor view is the algebraic view for Boolean
addresses and MLEs, while the flat-vector view is the implementation view for
storage, commitments, and linear-algebra kernels.

\section{The Multilinear Extension}

In this paper, MLE means \emph{multilinear extension}, not maximum likelihood
estimation.  It is one of the central objects in modern polynomial-commitment
and sumcheck protocols.

\begin{definition}[Multilinear polynomial]
A polynomial $P(X_0,\ldots,X_{d-1})$ is multilinear if every variable has degree
at most one.  For example,
\[
  3+5X_0+7X_1+11X_0X_1
\]
is multilinear in $X_0,X_1$, but $X_0^2+X_1$ is not.
\end{definition}

\begin{definition}[Multilinear extension]
Let $f\in\F^{\B^d}$ be an order-$d$ tensor.  Its multilinear extension
$\MLE(f):\F^d\to\F$ is the unique multilinear polynomial
that agrees with $f$ on the Boolean hypercube:
\[
  \MLE(f)(b)=f(b)
  \qquad
  \text{for every }b\in\B^d.
\]
\end{definition}

The explicit formula is
\[
  \MLE(f)(z)
  =
  \sum_{b\in\B^d}
    f(b)\chi_b(z),
\]
where
\[
  \chi_b(z)
  =
  \prod_{h=0}^{d-1}
  \bigl((1-b_h)(1-z_h)+b_hz_h\bigr).
\]
The polynomial $\chi_b$ is an indicator polynomial.  It equals $1$ at Boolean
point $b$ and equals $0$ at every other Boolean point.

If an implementation receives the entries as a flat vector of length $2^d$, the
algebraic object is still a tensor $m\in\F^{\B^d}$ obtained by deserializing
those entries into Boolean addresses.  Then
\[
  \MLE(m)(z)
  =
  \sum_{b\in\B^d} m[b]\mu_z[b],
\]
where $\mu_z[b]=\chi_b(z)$ is the public evaluation weight for Boolean address
$b$.  Thus an MLE evaluation is an inner product:
\[
  \MLE(m)(z)=\ip{\mu_z}{m}.
\]

\subsection{A Two-Variable Example}

Suppose $m\in\F^{\B^2}$, with entries
$m[0,0]=m_{00}$, $m[0,1]=m_{01}$, $m[1,0]=m_{10}$, and
$m[1,1]=m_{11}$.  Its MLE is
\[
\begin{aligned}
  \MLE(m)(X_0,X_1)
  &=
  m_{00}(1-X_0)(1-X_1)
  + m_{01}(1-X_0)X_1\\
  &\quad
  + m_{10}X_0(1-X_1)
  + m_{11}X_0X_1.
\end{aligned}
\]
At $(0,0)$ this equals $m_{00}$; at $(0,1)$ it equals $m_{01}$; and so on.
At a non-Boolean point, it is an interpolation of the tensor entries.

\begin{figure}[h]
\centering
\begin{tikzpicture}[
  x={(1.55cm,0cm)},
  y={(0.75cm,0.38cm)},
  z={(0cm,1.05cm)},
  >=Latex,
  line cap=round,
  line join=round,
  font=\small
]
  \def\hzerozero{0.25}
  \def\honezero{1.45}
  \def\hzeroone{0.85}
  \def\honeone{0.30}

  \draw[gray!45] (0,0,0) -- (1,0,0) -- (1,1,0) -- (0,1,0) -- cycle;
  \draw[-Latex,gray!70] (0,0,0) -- (1.18,0,0) node[below] {$x$};
  \draw[-Latex,gray!70] (0,0,0) -- (0,1.18,0) node[left] {$y$};
  \draw[-Latex,gray!70] (0,0,0) -- (0,0,1.65) node[left] {value};

  \foreach \x/\y/\h in {
    0/0/\hzerozero,
    1/0/\honezero,
    0/1/\hzeroone,
    1/1/\honeone
  } {
    \draw[gray!55,dashed] (\x,\y,0) -- (\x,\y,\h);
    \fill[black] (\x,\y,0) circle (0.9pt);
    \fill[red!75!black] (\x,\y,\h) circle (1.3pt);
  }
  \node[below left] at (0,0,\hzerozero) {$m_{00}$};
  \node[below right] at (1,0,\honezero) {$m_{10}$};
  \node[above left] at (0,1,\hzeroone) {$m_{01}$};
  \node[above right] at (1,1,\honeone) {$m_{11}$};

  \foreach \x in {0,0.25,0.5,0.75,1} {
    \draw[blue!55]
      plot[variable=\y,domain=0:1,samples=25]
      ({\x},{\y},{(1-\x)*(1-\y)*\hzerozero + \x*(1-\y)*\honezero
        + (1-\x)*\y*\hzeroone + \x*\y*\honeone});
  }
  \foreach \y in {0,0.25,0.5,0.75,1} {
    \draw[blue!55]
      plot[variable=\x,domain=0:1,samples=25]
      ({\x},{\y},{(1-\x)*(1-\y)*\hzerozero + \x*(1-\y)*\honezero
        + (1-\x)*\y*\hzeroone + \x*\y*\honeone});
  }
\end{tikzpicture}
\caption{A two-variable multilinear extension as a bilinear interpolating surface.  The tensor $m$ gives values at Boolean corners; $\MLE(m)$ fills in the unique surface linear in each coordinate.}
\label{fig:two-variable-mle}
\end{figure}

\subsection{Folding Is MLE Evaluation One Variable At A Time}

Folding is most naturally expressed on the tensor itself.  Throughout this
paper, when we write a fold formula, we use the convention that the folded
coordinate is the last tensor coordinate.  More general coordinate folds are the
same operation after reordering coordinates.

Let $x\in\F^{\B^d}$.  For $a\in\B^{d-1}$ and $b\in\B$,
write $(a\mid b)\in\B^d$ for the concatenation of the length-$(d-1)$ address
$a$ and the final bit $b$.

For example, if $d=3$ and $a=(1,0)$, then
\[
  (a\mid 0)=(1,0,0),
  \qquad
  (a\mid 1)=(1,0,1).
\]
Folding the last coordinate at a field element $\lambda$ produces a tensor
$x^\star\in\F^{\B^{d-1}}$:
\[
  x^\star[a]
  =
  (1-\lambda)x[(a\mid 0)]
  +
  \lambda x[(a\mid 1)].
\]
This operation evaluates one MLE variable at $\lambda$ and removes that tensor
coordinate.  If we fold all coordinates, using one challenge for each coordinate
in the order those coordinates are removed, the final scalar is
\[
  \MLE(x)(\lambda_0,\ldots,\lambda_{d-1}).
\]

For implementors, the same fold becomes pairwise interpolation in the serialized
array.  If the chosen coordinate is laid out so that entries with that bit equal
$0$ occupy the first half and entries with that bit equal $1$ occupy the second
half, then the coordinate fold is implemented by
\[
  x'[i]=(1-\lambda)x[i]+\lambda x[i+L],
  \qquad 0\le i<L.
\]
That flat formula is a storage layout for the tensor fold, not the definition
of the fold itself.
